\subsection{Realizing the nonsplit coefficient extension in an admissible theory}

\begin{theorem}[AT09: Proved; a nonsplit coefficient module with an explicit common source]
\label{AT09}
Fix integers $N>j\geq1$ and put $u=N-j$.
For each $c\in\mathbb A_f$ there is an admissible
simplicial strict PD differential $S$-algebra $A_c$
and an actual morphism
\[
 G\longrightarrow A_c,
 \tag{AT.10}
\]
where $G$ is the standard integral theory.
Every weight-form component in every simplex is finite
free abelian, and $A_c^{w,q}=0$ for $q>w$.
The coefficient map is a vertical quasi-isomorphism
in every weight and simplex, and $A_c^{w,q}$ is
simplicially acyclic for $w>q$.

In simplex degree zero, the form-zero algebra has an
$S$-submodule $F_1(A_{c,0}^{*,0})$ fitting into the
exact sequence
\[
 0\longrightarrow S\longrightarrow
 F_1(A_{c,0}^{*,0})\longrightarrow
 I_N\langle-j\rangle\longrightarrow0,
 \tag{AT.11}
\]
which is the shifted explicit extension
$E_{c,N,j}\langle-j\rangle$ of AT03.
For $c\notin\mathbb Q$, the coefficient inclusion
$S\to A_{c,0}^{*,0}$ has no weight-preserving
$S$-linear retraction. In particular there is no
strict coefficient-preserving morphism $A_c\to G$.
The morphism (AT.10) nevertheless supplies an
admissible common source for this pair.
\end{theorem}

\begin{proof}
We give the nonfree PD disk, the gluing map, and both
admissibility checks explicitly.

For a finite free abelian group $V$, put
\[
 \Omega(V)=\Gamma_{\mathbb Z}(V_x)
             \otimes\Lambda_{\mathbb Z}(V_{dx}),
 \qquad d x(v)=dx(v),\quad d(dx(v))=0.
\]
Its arity is the total number of $x$ and $dx$ factors,
counting divided powers with multiplicity; write
$\Omega_m(V)$ for its arity-$m$ complex.
Temporarily give $x(v)$ and $dx(v)$ weight $-j$.
Inside the resulting integer-weight-graded algebra
$S\otimes\Omega(V)$ take the submodule
\[
 C_{N,j}(V)=S\oplus
 \bigoplus_{m\geq1} I_{mN}\langle-mj\rangle
                      \otimes_{\mathbb Z}\Omega_m(V).
 \tag{AT.12}
\]
In this formula $\Omega_m(V)$ is assigned weight zero;
the shift records its temporary weight $-mj$.
Thus its displayed elements are precisely the linear
combinations of $T^{[a]}\omega$ with
$\omega\in\Omega_m(V)$ and $a\geq mN$.
Their actual weights are $a-mj\geq mu>0$.

The module (AT.12) is a PD differential subalgebra.
The differential preserves arity. A product of terms
of coefficient weights $a,b$ and arities $m,l$
has coefficient weight $a+b\geq(m+l)N$.
The divided power of index $r$ of a term of coefficient
weight $a$ and arity $m$ has coefficient weight $ra$
and arity $rm$; the inequality is again preserved.
The PD sum formula reduces arbitrary finite even sums
to these checks. All these operations are restrictions
of the standard operations of the ambient free PD
differential algebra, so their identities are retained.
Every actual weight is nonnegative, and a term of
form degree $q$ has $q\leq m\leq a-mj$.

Every positive-arity complex $\Omega_m(V)$ is
integrally contractible. To check this without dividing
by $m$, choose an ordered basis of $V$ and decompose
by the arity in each basis variable. In a component
of positive total arity, let $i$ be the first variable
with positive arity. Its one-variable complex is the
identity pair
\[
 \gamma_a(x_i)\longmapsto\gamma_{a-1}(x_i)dx_i.
\]
Apply its inverse as a degree-minus-one homotopy and
the identity on all other factors, with the tensor
sign rule. Its commutator with the full differential
is the identity: the other factor differentials
anticommute with this homotopy. This contracts each
multi-arity component and hence $\Omega_m(V)$.
Tensoring these homotopies with $I_{mN}$ gives an
$S$-linear, weight-preserving contraction of the
positive-arity part of (AT.12). Thus
$C_{N,j}(V)\simeq S$ vertically, before and after
any scalar extension used below.

Use $\Phi_n=\mathbb Z^{n+1}$, with faces deleting a
coordinate and degeneracies repeating one, and
$G=\Gamma(\Phi)\otimes\Lambda(\Phi/\mathbb Z1)$
with its standard differential and coefficient map.
These are the explicitly checked objects of A-L02.
For $R=\mathbb Z,\mathbb Q,\widehat{\mathbb Z},\mathbb A_f$
put
\[
 L_R=G_R\otimes_{S_R} C_{N,j}(\Phi)_R.
 \tag{AT.13}
\]
All parameters in (AT.13) use the same simplicial
index. The algebra $G_n$ is free over $S$ in each
form degree, by a complement to its marked coefficient
line. Therefore the indicated coefficient-tail
submodules embed in the corresponding free-variable
PD algebra over $G_R$. This specifies the tensor PD
structure in (AT.13), and all simplicial maps respect
it by functoriality in the parameter $\Phi$.

In the two rational coefficient rings, put $t=T$ and
replace the formal variable $x(v)$ by
$z(v)=t^N x(v)$. The inclusion of (AT.12) induces the
actual PD differential isomorphism
\[
 C_{N,j}(\Phi)_R
   \cong S_R\otimes\Gamma_R(\Phi_z)
                    \otimes\Lambda_R(\Phi_{dz}),
 \qquad R=\mathbb Q,\mathbb A_f,
 \tag{AT.14}
\]
where $z,dz$ have weight $u$ and form degrees $0,1$.
Indeed in arity $m$ the ideal $(I_{mN})_R$ is exactly
$t^{mN}R[t]$, and the change of variable in (AT.14)
is multiplication by that power on every monomial.
This also checks multiplicative and differential
compatibility directly.

We need a specific simplicial linear map
$f_u:\Phi\to G^{u,0}$. Write $e_0,\ldots,e_n$ for
the coordinate basis in simplex degree $n$, set
$v_i=e_i+\cdots+e_n$ and $v_{n+1}=0$, and define
\[
 f_u(e_i)=\gamma_u(v_i)-\gamma_u(v_{i+1}).
 \tag{AT.15}
\]
Deleting a coordinate sends the corresponding
difference to zero and the other differences to
their reindexed versions. Repeating coordinate $i$
splits its difference into two adjacent differences,
whose middle terms cancel. These are precisely the
face and degeneracy formulas for $\Phi$; hence
(AT.15) is simplicial. Its value on the marked
constant vector is $\gamma_u(T)=T^{[u]}$ by
telescoping, and in simplex degree zero this is its
value on $e_0$. Put $g_u=u!f_u$, so that
$(g_u)_0(e_0)=t^u$.

Using (AT.14), define the PD DGA automorphism
$\tau_c:L_{\mathbb A_f}\to L_{\mathbb A_f}$
by fixing $G_{\mathbb A_f}$ and setting
\[
 z(v)\longmapsto z(v)+c g_u(v),\qquad
 dz(v)\longmapsto dz(v)+c\,d_Gg_u(v).
 \tag{AT.16}
\]
The ring $\mathbb A_f$ is a rational algebra, so
these assignments extend through the ordinary free
graded-commutative algebra in (AT.14), preserving
its canonical divided powers. The differential
commutes with the assignments on both generators.
The inverse is $\tau_{-c}$. Naturality of (AT.15)
proves that (AT.16) is simplicial; both sides of
each assignment have the same weight and form
degree. Define
\[
 A_c=L_{\mathbb Q}
       \mathop{\times}_{L_{\mathbb A_f},\,\tau_c}
       L_{\widehat{\mathbb Z}}.
 \tag{AT.17}
\]
Explicitly its elements are pairs $(a,b)$ with
$b=\tau_c(a)$ after the specified coefficient
extensions. It is a strict PD DGA by its coordinate
operations. Its coefficients are the diagonal $S$;
its weight-zero algebra is $\mathbb Z$, and killing
positive weights is the required augmentation.
The automorphism fixes $G$, giving (AT.10).

Here are the exactness and finiteness details of this
pullback. Filter each fixed weight-form component
by polynomial arity in $z,dz$, using arity at most
$m$. The filtration has at most $\lfloor w/u\rfloor$
positive stages in weight $w$.
The translation (AT.16) preserves this filtration
and induces its identity on every associated graded
component. Each associated graded component of
$L_{\mathbb Z,n}^{w,q}$ is finite free abelian:
choose the finite-weight free $S$-basis of $G_n$
and the divided-power/exterior basis in (AT.12).
The associated graded of the difference map
\[
 L_{\mathbb Q,n}^{w,q}\oplus
 L_{\widehat{\mathbb Z},n}^{w,q}
 \longrightarrow L_{\mathbb A_f,n}^{w,q},
 \qquad(a,b)\longmapsto b-\tau_c(a)
 \tag{AT.18}
\]
is therefore a finite free module tensored with
the surjective arithmetic map of AT01. Induction
on this finite filtration proves surjectivity of
(AT.18). The same induction, applied to its kernels,
identifies
\[
 \operatorname{gr}_m A_c
      \cong\operatorname{gr}_m L_{\mathbb Z}
 \tag{AT.19}
\]
as simplicial bigraded modules. To see the kernel
surjectivity in this induction explicitly, lift an
associated-graded kernel pair; its difference lies
one filtration stage lower, where surjectivity
already proved permits a correction. Each weight
group in (AT.19) is finite free, so finite induction
also proves finite freeness of $A_{c,n}^{w,q}$.
The inequality $q\leq w$ is retained by (AT.17).

Vertically, the contractions of (AT.12) and the
standard $S$-linear contraction $G_n\simeq S$
show that $L_{R,n}\simeq S_R$ in every weight.
On this cohomology (AT.18) is the coefficient
arithmetic difference map, because $\tau_c$ fixes
coefficients. The long exact sequence of the
surjective cochain map (AT.18), together with AT01,
gives $H^0(A_{c,n}^{w,*})=S^w$ and
$H^q(A_{c,n}^{w,*})=0$ for every $q>0$.
All weights and all simplex degrees are included.

Horizontally, the arity-zero part of (AT.19) is $G$.
Every positive-arity part is a direct sum, in each
form degree, of factors containing
$\Gamma^{m-q'}(\Phi)\otimes\Lambda^{q'}(\Phi)$
with $m>0$. At least one of these two homogeneous
functors has positive degree and sends the additive
simplicial contraction of $\Phi$ to a contraction.
Tensoring its homotopy operators with degeneracies
on the other factors contracts the whole component.
Consequently every positive-arity row in (AT.19)
is contractible. The finite filtration in each
weight identifies the horizontal homology of
$A_c^{w,q}$ with that of $G^{w,q}$, compatibly with
(AT.10). The latter vanishes for $w>q$, by the
explicit standard-theory computation in A-L02.
This proves all admissibility assertions.

Finally in simplex degree zero $G_0=S$ and $\Phi_0$
has one generator. In arity at most one and form
degree zero, write a rational element as
$a+t^{-N}v\,z$, where $v\in(I_N)_{\mathbb Q}$
has its weight shifted by $-j$. Formula (AT.16)
changes its constant coefficient by
$c t^{-N}v\,t^u=c t^{-j}v$.
The equality of the two linear coefficients in
(AT.17) forces $v\in I_N$, by AT01 in each weight.
The equality of the constants is exactly
$b-a=c t^{-j}v$. Thus (AT.11) is the shifted
extension of AT03, with all three maps identified.
For $c\notin\mathbb Q$ it has no graded $S$-linear
retraction onto its coefficient copy of $S$.
Any such retraction from the whole $A_{c,0}^{*,0}$
would restrict to one on this submodule, a
contradiction. A strict map $A_c\to G$ would give
such a retraction in simplex degree zero since
$G_0=S$ in form zero. The map in the other direction
has already been explicitly constructed in (AT.10),
completing the proof.
\end{proof}
